\section{案例与运营：在 Berkeley 之外部署 GR00T}

\subsection{行业案例小结}

Jim Fan 分析了两个现实案例：

\begin{itemize}
    \item \textbf{半导体晶圆实验室}：使用 GR00T 模型配合自主抓取手臂，完成 wafer 交换。靠 built-in alignment rules 防止碰撞。
    \item \textbf{物流分拣中心}：目标在拥挤环境中拾取多种形状物体，借助 LLM 指令生成执行计划并通过 per-instance vision 进行误差检测。
\end{itemize}

\begin{warningbox}{部署时的真实尴尬}
真实场景里最常见的失败不是模型推理错误，而是 sensor calibration drift、机械 backlash 与工具更换导致的微小偏移。再强的模型也要有可靠的 fallback。
\end{warningbox}

\subsection{运营指标与监控}

为保障部署，把观测链分为三层：

\begin{itemize}
    \item \textbf{Trace}：每个动作/指令都附带 timestamp、token 序列、tool call log
    \item \textbf{Metrics}：成功率、cycle time、safety intervention frequency
    \item \textbf{Alerts}：预设 thresholds（碰撞、任务超时、planner oscillation）触发点亮红灯
\end{itemize}

\begin{knowledgebox}{Trace-driven postmortem}
保存 trace buffer 能在 incident 发生后快速回放：每个 token、工具调用、状态转移都可回溯，帮助工程师定位 ``在第 37 步哪行指令导致偏差''。
\end{knowledgebox}

\subsection{操作流程与 Playbook}

每个操控项目都遵循五步 Operational Playbook：

\begin{enumerate}
    \item \textbf{Plan}：定义任务、配置 CTA 文档、生成 failure hypotheses
    \item \textbf{Deploy}：写 live config、校准 sensors、注入 override switch
    \item \textbf{Observe}：通过 metrics + open telemetry 监控 drop rates、collisions、latency
    \item \textbf{Adjust}：如果触发 alert，进入 “pause + analyze” 模式，回放 trace，调整 instruction
    \item \textbf{Document}：每次 incident 写 postmortem，更新 Checklist
\end{enumerate}

\begin{warningbox}{Pause + Analyze 是最安全的制动}
遇到性能下降时，不要盲目调 large language predicate，只要按 checklist pause 机器人、回放 trace、查 lockstep metric，就能避免 cascade failure。
\end{warningbox}

\subsection{本章小结}

运营稳定需要 Case study 背书、监控/trace/backstop，同时配合 playbook 才能应对真实部署中的 incident。

\newpage

